% Generated draft: requires mathematical and source review.
\documentclass[11pt]{article}
\usepackage[margin=25mm]{geometry}
\usepackage{fontspec}
\setmainfont{Latin Modern Roman}
\usepackage{amsmath,amssymb,mathtools,unicode-math}
\setmathfont{Latin Modern Math}
\usepackage{enumitem,needspace,xurl,hyperref}
\setlength{\parindent}{0pt}
\setlength{\parskip}{5pt}
\setlength{\emergencystretch}{3em}
\newcommand{\N}{\mathbb{N}}
\newcommand{\Z}{\mathbb{Z}}
\newcommand{\Q}{\mathbb{Q}}
\newcommand{\R}{\mathbb{R}}
\newcommand{\C}{\mathbb{C}}
\newcommand{\F}{\mathbb{F}}
\newcommand{\A}{\mathbb{A}}
\newcommand{\PP}{\mathbb{P}}
\newcommand{\E}{\mathbb{E}}
\newcommand{\eps}{\varepsilon}
\newcommand{\dd}{\,\mathrm{d}}
\begin{document}
\section*{UnsolvedMath ID: 30005811}
\textbf{Source problem code:} catalogue entry
\subsection{Definitions and mathematical statement}
A simple pseudoline arrangement consists of pseudolines such that every pair intersects exactly once, transversely, and no three meet at one point. Its faces are the connected components of the complement of the pseudolines in the ambient surface. The supplied record does not specify the ambient surface, the equivalence relation on arrangements, or the conventions for counting sides of unbounded faces. These choices must be fixed before the requested probability is uniquely defined.

Conditional on those choices, let \(\mathcal A_n\) be the resulting finite, nonempty set of combinatorial arrangements of \(n\) pseudolines. For \(A\in\mathcal A_n\), let \(F(A)\) be its finite set of faces, and let \(s_A(f)\) denote the number of sides of \(f\), counted according to the chosen boundary convention. Choose \(A_n\) uniformly from \(\mathcal A_n\), and then, conditional on \(A_n=A\), choose \(U_n\) uniformly from \(F(A)\). Thus, for each fixed integer \(k\geq 3\), set
\[
q_{n,k}:=\Pr\bigl(s_{A_n}(U_n)=k\bigr)
=\frac{1}{|\mathcal A_n|}\sum_{A\in\mathcal A_n}
\frac{|\{f\in F(A):s_A(f)=k\}|}{|F(A)|}.
\]
The reconstructed statement in the supplied registry asks to prove that, for every fixed integer \(k\geq 3\), the limit
\[
p_k:=\lim_{n\to\infty}q_{n,k}
\]
exists and to determine its value, beginning with \(p_3\), the limiting triangle probability. This formulation remains conditional on specifying \(\mathcal A_n\) and the face conventions; no particular labeled, unlabeled, affine, or projective model is selected here.

The original excerpt explicitly asks only for \(p_3\) and asserts that the answer is also the limiting probability that the convex hull of a random abstract order type has size three. The record does not define the sampling model for abstract order types or the correspondence supporting this assertion, so it is retained as a source qualification rather than a separately precise equivalent formulation. The request concerning all \(p_k\) comes from the supplied reconstructed statement.
\subsection{Short English statement}
Choose a simple pseudoline arrangement uniformly, then choose one of its faces uniformly. The supplied reconstructed problem asks whether the probability of a \(k\)-sided face converges for every fixed \(k\geq 3\), and asks for the limits, starting with triangles. The arrangement model and face conventions require clarification. The original excerpt explicitly asks the triangle question and also asserts an unspecified connection with convex hulls of random abstract order types.
\subsection{Sources}
\begin{itemize}
\item Local UnsolvedMath registry record, permanent ID \(30005811\), source code \texttt{catalogue entry}, title \emph{Face-Size Distribution in Random Pseudoline Arrangements}; supplied as the complete source record for this draft. The registry labels the problem open and its reconstructed statement \texttt{corrected\_verified}, and reports that the statement was checked against the official report. These are supplied registry assertions, not independent verification in this draft. The proposer is unspecified; the supplied proposed year is 2024.
\item \emph{Discrete Geometry} (2024), Oberwolfach Reports, pp.\ 137--202, DOI 10.4171/owr/2024/3: \url{https://doi.org/10.4171/owr/2024/3}. This is the supplied report citation, source URL, and literature-source link. No more precise location within the report is supplied for the target problem.
\item Literature assessment dated exactly \texttt{2026-08-22}: the supplied curation labels the problem open and attributes to the report the statement that even convergence of the random-face-size probabilities is unproved, hence that the limiting triangle probability remains open. This is supplied dated curation, not independently verified literature assessment; the record itself requests independent verification of time-sensitive status and cited claims.
\item The original excerpt also contains unrelated material attributed to Sergey Avvakumov concerning grid graphs and thick drawings, and to Karim Adiprasito concerning Moore graphs. Its incidental reference is Larry Guth, \emph{The Width-Volume Inequality}, Geometric and Functional Analysis 17 (2007), pp.\ 1139--1179. That reference belongs to the grid-graph passage and is not supplied as evidence for the pseudoline problem; neither named problem attribution is transferred to the target problem.
\end{itemize}
\end{document}
